\section{Analysis Methods}

The number of neutrino interactions in the fiducial volume of the P$\emptyset$D, $N_{signal}$,
can be expressed as the product of the signal cross section per target, $\sigma$, the number of targets, $N_{targets}$, and the integrated flux, $\Phi$, of incident
neutrinos per unit area, as
\begin{equation}
N_{signal}=\sigma N_{targets}\Phi.\label{eq:Analysis:1-1}
\end{equation}
Hence the cross section becomes
\begin{equation}
\sigma=\frac{N_{signal}}{\Phi N_{targets}}.\label{eq:Analysis:2-1}
\end{equation}
Using our event selection on data, we obtain a candidate signal event sample in our fiducial volume.
This process is not 100\% efficient and also some non-signal (background) events 
are included.
To account for this, the MC simulation is used to estimate in our sample
the number of background events and the number of
signal events.  The backgrounds from the FHC (RHC) beam samples
include non-CCINC events from the neutrino (antineutrino) beam as well as
events created from the antineutrino (neutrino) flux.  
In addition, the MC simulation generates
total number of signal events that were produced. 
If the rate of 
restricted phase space
selected data events is $N_{selected}^{data}$
and the predicted number of selected background events is $B^{MC}$,
the observed number of signal 
candidates 
in our fiducial volume is
\begin{equation}
N_{selected\ signal}=N_{selected}^{data}-B^{MC}\label{eq:Analysis:3-1}
\end{equation}
which include migration events.
%If $N^{MC}_{generated\ signal}$ is the number of true signal events generated
%to have occurred 
%n the generated kinematic phase space fiducial volume 
%in our fiducial volume 
%and  $N_{selected\ signal}^{MC}$
%is the number of reconstructed restricted phase space MC event candidates, then
Next we redefine the selection efficiency
$\epsilon$ as
\begin{equation}
\epsilon=\frac{N_{selected\ signal}^{MC}}{N_{generated\ signal}^{MC}}.\label{eq:Analysis:4-1}
\end{equation}
where the $N^{MC}_{selected\ signal}$ is the number of signal candidates whose reconstructed kinematics 
are in the restricted phase space and  $N^{MC}_{generated\ signal}$ is the total number of generated 
signal events whose true kinematics are in the restricted phase space. We note that
 $N^{MC}_{selected\ signal}$ includes a small fraction of migration events
as described at the end of Section IV.B.
With these definitions, the 
restricted phase space
signal event rate is
\begin{equation}
N_{signal} = \frac{N_{selected}^{data}-B^{MC}}{\epsilon}.\label{eq:Analysis:5-1}
\end{equation}
In Eqn.(\ref{eq:Analysis:5-1}) the numerator
is the number of signal candidates whose reconstructed kinematics are in the restricted 
phase space, and this is combined with the
denominator  $\epsilon$ from Eqn.(\ref{eq:Analysis:4-1})
to give the proper estimate 
of $N_{signal}$ that represents the number of signal events 
whose kinematics are in the true restricted phase space.
The neutrino cross section is
\begin{equation}
\sigma (\nu_{\mu})=\frac{N_{selected}^{data}-B^{MC}}{\epsilon N_{targets}\Phi}.\label{eq:Analysis:6-1}
\end{equation}

In addition to the cross sections given above, the measured ratio of cross sections $R(\nu,\overline{\nu})$ and rates $r(\nu,\overline{\nu})$
are defined as
\begin{equation}
R (\nu,\overline{\nu}) \equiv \frac{\sigma(\overline{\nu}_{\mu})}{\sigma(\nu_{\mu})}=\frac{\overline{N}_{selected}^{data}-\overline{B}^{MC}}{N_{selected}^{data}-B^{MC}}\times\frac{\overline{\epsilon}}{\epsilon}\times\frac{\varPhi}{\overline{\varPhi}}\label{eq:Analysis:7-1}
\end{equation}
and
\begin{equation}
r (\nu,\overline{\nu}) \equiv \frac{n(\overline{\nu}_{\mu})}{n(\nu_{\mu})}=\frac{\overline{N}_{selected}^{data}-\overline{B}^{MC}}{N_{selected}^{data}-B^{MC}}\times\frac{\overline{\epsilon}}{\epsilon}.\label{eq:Analysis:7-2}
\end{equation}
The overlined quantities are obtained from the antineutrino
selections as described above and those without overlines represent the
neutrino mode selection.  Finally, other observables are introduced and defined; 
the sum $\Sigma(\nu,\overline{\nu})$, 
difference $\Delta(\nu,\overline{\nu})$, 
and asymmetry $A(\nu,\overline{\nu})$
formed from the $\nu_\mu$ and \nubar 
cross sections, as
\begin{equation}
\Sigma(\nu,\overline{\nu}) \equiv \sigma(\nu_{\mu})+\sigma(\overline{\nu}_{\mu}),\label{eq:Analysis:sum}
\end{equation}
\begin{equation}
\Delta(\nu,\overline{\nu}) \equiv \sigma(\nu_{\mu})-\sigma(\overline{\nu}_{\mu})\label{eq:Analysis:diff}
\end{equation}
and
\begin{equation}
A(\nu,\overline{\nu}) \equiv 
\frac{\sigma(\nu_{\mu})-\sigma(\overline{\nu}_{\mu})}{\sigma(\nu_{\mu})+\sigma(\overline{\nu}_{\mu})}.\label{eq:Analysis:asym}
\end{equation}